

\subsubsection{Disco en barra rotatoria vertical}

Un disco circular de radio $R$ está unido al extremo de una barra rígida de longitud $L$. La barra gira libremente con una velocidad angular constante $\dot{\phi}$ alrededor de un eje vertical que pasa por el punto $A$. A su vez, el disco puede girar respecto a su propio centro $O$ sobre un eje ortogonal a la barra. La posición de un punto $P$ ubicado en el borde del disco se describe mediante el ángulo $\theta$, medido desde una referencia fija en el disco como se muestra en la Figura \ref{fig:disco_barra_rotatoria}.

Considerando que se conocen las magnitudes $\dot{\phi}$, $\theta$ y $\ddot{\theta}$, determine lo siguiente para un observador fijo al suelo:

\begin{enumerate}
    \item Defina los marcos de referencia y los sistemas de coordenadas adecuados para describir el movimiento compuesto del sistema.
    \item Determine el vector posición del punto $P$ respecto al punto $A$.
    \item Obtenga los vectores de velocidad angular y aceleración angular de la barra y del disco.
    \item Calcule los vectores de velocidad y aceleración del punto $P$.
\end{enumerate}

\begin{figure}[!ht]
    \centering
    \begin{tikzpicture}[>=Latex, scale=1.5]
        % --- Parámetros Geométricos ---
        \def\L{3.5}      % Longitud de la barra
        \def\R{1.2}      % Radio del disco
        \def\thetaAng{40} % Ángulo theta de P

        % --- Coordenadas principales ---
        \coordinate (A) at (0,0);
        \coordinate (O) at (\L,0);
        \coordinate (P) at ($ (O) + ({90-\thetaAng}:\R) $);

        % --- 1. Soporte y Eje en A ---
        % Sombreado de la base
        \fill[pattern=north east lines] (-0.4, 0.3) rectangle (-0.1, 0.7);
        \fill[pattern=north west lines] (0.1, 0.3) rectangle (0.4, 0.7);
        \fill[pattern=north east lines] (-0.4, -0.3) rectangle (-0.1, -0.7);
        \fill[pattern=north west lines] (0.1, -0.3) rectangle (0.4, -0.7);
        
        \draw[thick] (-0.4, 0.3) -- (-0.1, 0.3) -- (-0.1, 0.7);
        \draw[thick] (0.4, 0.3) -- (0.1, 0.3) -- (0.1, 0.7);
        \draw[thick] (-0.4, -0.3) -- (-0.1, -0.3) -- (-0.1, -0.7);
        \draw[thick] (0.4, -0.3) -- (0.1, -0.3) -- (0.1, -0.7);

        % Eje vertical
        \shadedraw[left color=gray!10, right color=gray!40, draw=black, rounded corners=2pt] (-0.1, 1.2) -- (0.1, 1.2)--(0.1,0.27)--(0.4,0.27)--(0.4,-0.27)--(0.1,-0.27)--(0.1, -1.2)--(-0.1, -1.2)--(-0.1,-0.27)--(-0.4,-0.27)--(-0.4,0.27)--(-0.1,0.27)--cycle;
        \draw[dashed, thin] (0, 1.2) -- (0, -1.2);
        \draw[->, thick] (-0.2, 0.9) arc (120:420:0.4 and 0.15);
        \node[right] at (0.3, 1.0) {$\dot{\phi} = \text{cte}$};
    % --- 3. El Disco (R) ---
        \shadedraw[inner color=gray!5, outer color=gray!30, draw=black, thick] (O) circle (\R);
        \draw[dashed] (O) -- ($(O) + (0, \R)$); % Línea vertical de referencia en disco
        % --- 2. La Barra (L) ---
        \shadedraw[left color=gray!10, right color=gray!40, draw=black, rounded corners=2pt] 
            (-0.15, -0.12) rectangle (\L+0.1, 0.12);
        
        % Dimensión L
        \draw[<->] (0, -0.8) -- (\L, -0.8) node[midway, fill=white] {$L$};
        \draw[dotted] (0, -0.15) -- (0, -0.9);
        \draw[dotted] (\L, -0.15) -- (\L, -0.9);

       
        
        % Radio R y Ángulo Theta
        \draw[->] (O) -- (P) node[midway, right, xshift=2pt] {$R$};
        \draw[->] ($(O) + (0, 0.6)$) arc (90:{90-\thetaAng}:0.6);
        \node at ($ (O) + (72:0.8) $) {$\theta$};

        % --- 4. Puntos y Etiquetas ---
        \filldraw (A) circle (1.5pt) node[below right, xshift=2pt] {$A$};
        \filldraw (O) circle (1.5pt) node[right] {$O$};
        \filldraw (P) circle (1.8pt) node[right] {$P$};

    \end{tikzpicture}
    \caption{Esquema del disco circular de radio $R$ montado sobre una barra rotatoria de longitud $L$.}
    \label{fig:disco_barra_rotatoria}
\end{figure}

\ifmostrarSolucion
\begin{tcolorbox}[% Borde gris oscuro
        title= Solucion Disco en barra rotatoria vertical, breakable]

\textbf{Marcos de referencia y Sistemas de Coordenadas}

Al igual que en siempre, partimos por los marcos de referencia

\begin{itemize}
    \item $\mathcal{N}$ El suelo $\longrightarrow\{\vec{e}_x;\vec{e}_y;\vec{e}_z\}$
    \item $\mathcal{B}$ Brazo OA $\longrightarrow\{\vec{E}_r;\vec{E}_\phi;\vec{E}_z\}$
    \item $\mathcal{D}$ Disco $\longrightarrow\{\vec{e}_r;\vec{e}_\theta;\vec{e}_z\}$
\end{itemize}
\vspace{1cm}
Una vez decidido como hago los vectores, los dibujo y hago la asociacion de cada uno (recordar el que este asociado no significa que solo lo pueda usar en ese cuerpo)
\begin{center}
    \begin{tikzpicture}[>=Latex, scale=1.5]
        % --- Parámetros Geométricos ---
        \def\L{3.5}      % Longitud de la barra
        \def\R{1.2}      % Radio del disco
        \def\thetaAng{40} % Ángulo theta de P

        % --- Coordenadas principales ---
        \coordinate (A) at (0,0);
        \coordinate (O) at (\L,0);
        \coordinate (P) at ($ (O) + ({90-\thetaAng}:\R) $);

        % --- 1. Soporte y Eje en A ---
        % Sombreado de la base
        \fill[pattern=north east lines] (-0.4, 0.3) rectangle (-0.1, 0.7);
        \fill[pattern=north west lines] (0.1, 0.3) rectangle (0.4, 0.7);
        \fill[pattern=north east lines] (-0.4, -0.3) rectangle (-0.1, -0.7);
        \fill[pattern=north west lines] (0.1, -0.3) rectangle (0.4, -0.7);
        
        \draw[thick] (-0.4, 0.3) -- (-0.1, 0.3) -- (-0.1, 0.7);
        \draw[thick] (0.4, 0.3) -- (0.1, 0.3) -- (0.1, 0.7);
        \draw[thick] (-0.4, -0.3) -- (-0.1, -0.3) -- (-0.1, -0.7);
        \draw[thick] (0.4, -0.3) -- (0.1, -0.3) -- (0.1, -0.7);

        % Eje vertical
        \shadedraw[left color=gray!10, right color=gray!40, draw=black, rounded corners=2pt] (-0.1, 1.2) -- (0.1, 1.2)--(0.1,0.27)--(0.4,0.27)--(0.4,-0.27)--(0.1,-0.27)--(0.1, -1.2)--(-0.1, -1.2)--(-0.1,-0.27)--(-0.4,-0.27)--(-0.4,0.27)--(-0.1,0.27)--cycle;
        \draw[dashed, thin] (0, 1.2) -- (0, -1.2);
        \draw[->, thick] (-0.2, 0.9) arc (120:420:0.4 and 0.15);
        \node[right] at (0.3, 1.0) {$\dot{\phi} = \text{cte}$};
    % --- 3. El Disco (R) ---
        \shadedraw[inner color=gray!5, outer color=gray!30, draw=black, thick] (O) circle (\R);
        \draw[dashed] (O) -- ($(O) + (0, \R)$); % Línea vertical de referencia en disco
        % --- 2. La Barra (L) ---
        \shadedraw[left color=gray!10, right color=gray!40, draw=black, rounded corners=2pt] 
            (-0.15, -0.12) rectangle (\L+0.1, 0.12);
        
        % Dimensión L
        \draw[<->] (0, -0.8) -- (\L, -0.8) node[midway, fill=white] {$L$};
        \draw[dotted] (0, -0.15) -- (0, -0.9);
        \draw[dotted] (\L, -0.15) -- (\L, -0.9);

       
        
        % Radio R y Ángulo Theta
        \draw[->] (O) -- (P) node[midway, right, xshift=2pt] {$R$};
        \draw[->] ($(O) + (0, 0.6)$) arc (90:{90-\thetaAng}:0.6);
        \node at ($ (O) + (72:0.8) $) {$\theta$};

        % --- 4. Puntos y Etiquetas ---
        \filldraw (A) circle (1.5pt) node[below right, xshift=2pt] {$A$};
        \filldraw (O) circle (1.5pt) node[right] {$O$};
        \filldraw (P) circle (1.8pt) node[right] {$P$};

        %Aqui van los vectores
        
        

    \end{tikzpicture}
\end{center}

Nos falta lo más importante del paso 1: relacionar cada base con la anterior en la cadena. El
brazo $\mathcal{B}$ gira en torno al eje vertical (compartido con el suelo, $\vec{E}_z=\vec{e}_z$)
con ángulo $\phi$, así que
\begin{equation*}
\begin{split}
    \vec{E}_r&=\cos(\phi)\vec{e}_x+\sin(\phi)\vec{e}_y\\
    \vec{E}_\phi&=-\sin(\phi)\vec{e}_x+\cos(\phi)\vec{e}_y\\
    \vec{E}_z&=\vec{e}_z
\end{split}
\end{equation*}

El disco, en cambio, gira respecto al brazo en torno al eje $\vec{E}_\phi$ —el eje ortogonal a la
barra— con ángulo $\theta$ medido desde $\vec{E}_z$ (la línea punteada vertical de la figura), por
lo que
\begin{equation*}
\begin{split}
    \vec{e}_r&=\sin(\theta)\vec{E}_r+\cos(\theta)\vec{E}_z\\
    \vec{e}_\theta&=\cos(\theta)\vec{E}_r-\sin(\theta)\vec{E}_z\\
    \vec{e}_z&=\vec{E}_\phi
\end{split}
\end{equation*}

Vale la pena detenerse en la última línea: el $\vec{e}_z$ del disco —su propio eje de giro— no es
la vertical, sino que coincide con $\vec{E}_\phi$. Esto es justamente lo que dice el enunciado:
el disco gira sobre un eje ortogonal a la barra, no sobre su eje vertical.

\subsubsection*{Vector posición}

\textbf{Posicion}\\
Anclamos en $A$, que está quieto en el marco inercial, y sumamos hasta llegar a $P$ pasando por
$O$:
$$\vec{r}_{P/A}=\vec{r}_{O/A}+\vec{r}_{P/O}=L\vec{E}_r+R\vec{e}_r$$
Conviene dejarlo todo expresado en la base del brazo $\mathcal{B}$, reemplazando $\vec{e}_r$ por
la transformación de arriba:
$$\vec{r}_{P/A}=L\vec{E}_r+R\left(\sin(\theta)\vec{E}_r+\cos(\theta)\vec{E}_z\right)=(L+R\sin(\theta))\,\vec{E}_r+R\cos(\theta)\,\vec{E}_z$$

\subsubsection*{Velocidades y aceleraciones angulares}

Antes de derivar cualquier vector conviene tener a mano las velocidades angulares de la cadena,
ya que van a aparecer una y otra vez en el transporte.

La velocidad angular del brazo respecto al suelo es directa, pues $\phi$ se mide justo en torno a
$\vec{e}_z=\vec{E}_z$:
$$\velang{N}{B}=\dot\phi\,\vec{E}_z$$
Como $\dot\phi$ es constante y $\vec{E}_z$ nunca cambia (coincide siempre con la vertical fija),
la aceleración angular del brazo respecto al suelo es nula:
$$\accang{N}{B}=\vec{0}$$

El disco gira respecto al brazo en torno a $\vec{E}_\phi$:
$$\velang{B}{D}=\dot\theta\,\vec{E}_\phi$$
Como las velocidades angulares se suman a lo largo de la cadena de marcos, la velocidad angular
del disco respecto al suelo es
$$\velang{N}{D}=\velang{N}{B}+\velang{B}{D}=\dot\phi\,\vec{E}_z+\dot\theta\,\vec{E}_\phi$$

Para la aceleración angular del disco respecto al suelo hay que derivar $\velang{N}{D}$ vista
desde $\mathcal{N}$; como está expresada en vectores de $\mathcal{B}$, transportamos hacia
$\mathcal{B}$:
$$\accang{N}{D}=\der{}{N}\velang{N}{D}=\der{}{B}\velang{N}{D}+\velang{N}{B}\times\velang{N}{D}$$
El primer término es directo, ya que $\dot\phi$ es constante y $\vec{E}_z,\vec{E}_\phi$ no cambian
vistos desde $\mathcal{B}$:
$$\der{}{B}\velang{N}{D}=\ddot\theta\,\vec{E}_\phi$$
El segundo término es el producto cruz
$$\velang{N}{B}\times\velang{N}{D}=\dot\phi\vec{E}_z\times\left(\dot\phi\vec{E}_z+\dot\theta\vec{E}_\phi\right)=\dot\phi\dot\theta\left(\vec{E}_z\times\vec{E}_\phi\right)=-\dot\phi\dot\theta\,\vec{E}_r$$
usando que $\{\vec{E}_r,\vec{E}_\phi,\vec{E}_z\}$ es una base derecha ($\vec{E}_z\times\vec{E}_r=\vec{E}_\phi$, y por lo tanto $\vec{E}_z\times\vec{E}_\phi=-\vec{E}_r$). Sumando ambos términos:
$$\accang{N}{D}=-\dot\phi\dot\theta\,\vec{E}_r+\ddot\theta\,\vec{E}_\phi$$

\subsubsection*{Cinemática del punto $P$}

\textbf{Velocidad}\\
La velocidad inercial de $P$ se obtiene derivando la posición. Como $\vec{r}_{P/A}$ está escrito
en la base de $\mathcal{B}$, no podemos derivar directo y aplicamos el teorema del transporte
entre $\mathcal{N}$ y $\mathcal{B}$:
$$\marc{\vec{v}_P}{N}=\der{\vec{r}_{P/A}}{N}=\transporta{\vec{r}_{P/A}}{N}{B}$$
con $\velang{N}{B}=\dot\phi\,\vec{E}_z$ ya obtenido, y la derivada vista desde $\mathcal{B}$ —donde
$\vec{E}_r,\vec{E}_z$ son constantes, así que basta derivar las componentes—:
$$\der{\vec{r}_{P/A}}{B}=R\dot\theta\cos(\theta)\,\vec{E}_r-R\dot\theta\sin(\theta)\,\vec{E}_z$$
El producto cruz que falta es
$$\dot\phi\vec{E}_z\times\left[(L+R\sin(\theta))\vec{E}_r+R\cos(\theta)\vec{E}_z\right]=\dot\phi(L+R\sin(\theta))\,\vec{E}_\phi$$
Sumando ambos términos obtenemos la velocidad de $P$ respecto al suelo:
$$\marc{\vec{v}_P}{N}=R\dot\theta\cos(\theta)\,\vec{E}_r+\dot\phi(L+R\sin(\theta))\,\vec{E}_\phi-R\dot\theta\sin(\theta)\,\vec{E}_z$$

\textbf{Aceleración}\\
Repetimos el procedimiento sobre la velocidad recién obtenida, transportando otra vez entre
$\mathcal{N}$ y $\mathcal{B}$:
$$\marc{\vec{a}_P}{N}=\transporta{\marc{\vec{v}_P}{N}{}}{N}{B}$$
La derivada vista desde $\mathcal{B}$ sale de derivar cada componente de $\marc{\vec{v}_P}{N}{}$
(recordando que $\dot\phi$ es constante):
$$\der{}{B}\marc{\vec{v}_P}{N}{}=\left(R\ddot\theta\cos(\theta)-R\dot\theta^2\sin(\theta)\right)\vec{E}_r+\dot\phi R\dot\theta\cos(\theta)\,\vec{E}_\phi+\left(-R\ddot\theta\sin(\theta)-R\dot\theta^2\cos(\theta)\right)\vec{E}_z$$
y el producto cruz que falta es
$$\dot\phi\vec{E}_z\times\marc{\vec{v}_P}{N}{}=-\dot\phi^2(L+R\sin(\theta))\,\vec{E}_r+\dot\phi R\dot\theta\cos(\theta)\,\vec{E}_\phi$$
Sumando ambos términos, la aceleración de $P$ respecto al suelo es
$$\marc{\vec{a}_P}{N}=\Big[R\ddot\theta\cos(\theta)-R\dot\theta^2\sin(\theta)-\dot\phi^2(L+R\sin(\theta))\Big]\vec{E}_r+2\dot\phi R\dot\theta\cos(\theta)\,\vec{E}_\phi+\Big[-R\ddot\theta\sin(\theta)-R\dot\theta^2\cos(\theta)\Big]\vec{E}_z$$

Vale la pena mirar de dónde viene cada pieza: el término $-\dot\phi^2(L+R\sin\theta)$ en
$\vec{E}_r$ es el centrípeto de girar a distancia $L+R\sin\theta$ del eje vertical; el término en
$\vec{E}_\phi$, $2\dot\phi R\dot\theta\cos\theta$, es el de Coriolis y trae su factor 2 como
corresponde; y los términos restantes en $\vec{E}_r,\vec{E}_z$ son exactamente la aceleración
tangencial y centrípeta del movimiento circular de $P$ en torno a $O$ dentro del disco.

\end{tcolorbox}
\fi